\section{Related Work}
\label{sec:related-work}

    \subsection{Remote Control Systems for Automated Railways}
        As railway systems advance toward higher levels of automation (GoA 3 and 4), the role of remote control as a safety-critical fallback mechanism has gained significant attention. Cogan-Milius et al. \cite{CoganMilius2023} conducted a comprehensive study investigating the needs and preferences of future operators for remote control systems in automated trains. Their research highlights that while full automation aims to reduce human involvement, remote control capabilities remain essential for handling edge cases, system degradation, and emergency situations where human judgment and intervention are necessary. The study emphasizes that operators require intuitive interfaces, reliable communication channels, and situational awareness tools to effectively monitor and control trains remotely when automation systems fail or encounter unprecedented scenarios. These findings underscore the importance of designing remote control systems that not only provide technical capabilities but also align with human factors and operational workflows—a consideration that directly influences the architectural and interface design decisions in modern remote control platforms.

    \subsection{QUIC Protocol Performance in Real-Time Applications}
        The QUIC protocol, standardized as the foundation for HTTP/3, has emerged as a promising alternative to traditional TCP- and UDP-based communication for latency-sensitive applications. A comprehensive performance evaluation conducted in 2024 \cite{10.1145/3696355.3699698} analyzed QUIC's behavior in real-time network environments, demonstrating that its integration of transport and cryptographic handshakes significantly reduces connection setup time compared to TCP+TLS, achieving near-zero round-trip time (0-RTT) connection resumption in many scenarios. Furthermore, QUIC's built-in congestion control mechanisms and multiplexed stream architecture effectively mitigate head-of-line blocking issues that affect TCP-based protocols, particularly under lossy network conditions. Complementary research on Extended Reality (XR) streaming services has further validated QUIC’s suitability for real-time media transmission. Mejías et al. \cite{mejías2025streamingremoterenderingservices} compared Media over QUIC (MoQ), RTP over QUIC (RoQ), and WebRTC in remote rendering scenarios over Wi-Fi and 5G networks, reporting up to 30\% latency reduction and 60\% faster connection startup for QUIC-based protocols relative to WebRTC. Their results highlight QUIC’s advantages in connection establishment and adaptability across heterogeneous network environments. Together, these studies provide empirical evidence of QUIC’s growing maturity and potential as a transport solution for both interactive XR applications and safety-critical control systems requiring low latency, reliability, and seamless mobility support.


    \subsection{WebTransport for game streaming}
        Building upon QUIC and HTTP/3, the emerging WebTransport protocol introduces a flexible framework for multiplexed, bidirectional, and low-latency communication between clients and servers. Nguyen et al. \cite{10.1145/3744725.3744726} demonstrated the feasibility of employing WebTransport in a real-time game streaming system, where frames rendered in the Unity engine were encoded via FFmpeg and transmitted to web clients using unidirectional streams, while user input was relayed back through datagrams. Their implementation achieved an end-to-end latency of approximately 115 ms in local testing, confirming WebTransport’s capability to support interactive, time-critical applications comparable to WebRTC-based solutions. Compared to WebRTC, WebTransport offers a simpler server configuration by eliminating the need for auxiliary protocols such as ICE, SCTP, and DTLS, while providing fine-grained control over reliable and unreliable data transmission. The study also highlighted WebTransport’s scalability and its potential to reduce signaling overhead, which are particularly advantageous for distributed systems that require both media streaming and control signaling under varying network conditions. These attributes suggest that WebTransport can serve as a foundational communication layer for remote control and monitoring systems in automated railway operations, where deterministic latency, robustness, and efficient bidirectional messaging are critical for ensuring safety and responsiveness.

    \subsection{Remote Driving with Real-Time Video Streaming}
        The challenges of remote vehicle operation over wireless networks share fundamental similarities with remote train control, particularly regarding the critical interplay between video streaming latency, network conditions, and operator response times. Yu and Lee \cite{Yu2022RemoteDriving} developed and evaluated a comprehensive remote driving system that integrates real-time video feedback with low-latency control commands over wireless networks, addressing scenarios where autonomous systems require human intervention or where full autonomy is not yet achievable. Their system, implemented using the Robot Operating System (ROS), demonstrated that round-trip time (RTT)—encompassing video encoding, network transmission, decoding, operator reaction, and control command propagation—is the primary determinant of remote driving performance and safety. Through extensive testing over WiFi and LTE networks, the study revealed that maintaining RTT below 200 ms is essential for stable remote operation, with network jitter and packet loss significantly degrading operator situational awareness and control precision. The research further emphasized the importance of adaptive video encoding strategies, where bitrate and frame rate must dynamically adjust to network conditions to minimize feedback delay while preserving visual fidelity sufficient for safe navigation. These findings directly inform the design requirements for remote railway control systems, where operators must similarly process real-time video streams from onboard cameras while issuing time-critical control commands, and where wireless network variability poses comparable challenges to maintaining the low-latency, high-reliability communication necessary for safe remote intervention in automated train operations.
